% !TeX root = ../thuthesis-example.tex

\begin{comments}
% \begin{comments}[name = {指导小组评语}]
% \begin{comments}[name = {Comments from Thesis Supervisor}]
% \begin{comments}[name = {Comments from Thesis Supervision Committee}]
    论文以云控系统背景下的半挂式智能网联液罐车为研究对象，开展了车云协同半挂式液罐车防侧翻行驶规划与控制方法研究，选题具有理论意义与应用价值。

    论文建立了车云协同液罐车防侧翻场景任务体系架构，提出了面向交通交互场景的图扩散云端预测性决策与规划方法，并提出了面向液罐车横纵向耦合晃动特性的抑晃防侧翻轨迹跟踪控制方法从体系角度保障了液罐车行驶安全。针对液罐车防侧翻研究中长期存在的“架构层缺位、方法层孤立”问题，进行了战略、运行、服务、资源四阶段一体化建模，并通过仿真验证完成架构闭环，为后续算法模块的问题边界、接口衔接、时间尺度等参数划分及工程部署提供了可追溯的分析依据。针对液罐车在高速场景下面临的周车交互不确定性强、长时域风险演化建模难、现有表征方法信息冗余大等问题，提出了交通图表征方法、自回归图扩散交互式预测模型及强化学习长时域决策模型，建立了从周车未来行为预测到液罐车横纵向耦合决策再到防侧翻平滑参考轨迹生成的完整方法链，实现了面向诱发性风险的前瞻规避。针对液罐车车液耦合动力学计算复杂及现有研究仅考虑横向动力学的问题，构建了高精度反映罐液耦合动态的滑动伸缩天平快速计算模型及其残差强化学习重线性化方法，并基于此设计了多约束模型预测控制器以抑制状态风险进一步演化，形成了与长时域前瞻规划协同的短时域高频安全保障。并通过仿真与缩比液罐车试验，验证了所提方法的有效性。

    论文主要创新点包括：（1）创建了车云协同液罐车防侧翻场景任务体系架构；（2）提出了面向交通交互场景的图扩散云端预测性决策与规划方法；（3）提出了面向液罐车横纵向耦合晃动特性的抑晃防侧翻轨迹跟踪控制方法。

    戚笑景独立撰写的硕士学位论文，表明其在此领域具备扎实的理论基础，其提出的理论、技术和方法，证明其在此领域具备独立开展研究工作、解决实际问题的能力，其工作量达到了硕士学位论文的要求，同意其申请硕士学位论文答辩。

\end{comments}
